\documentclass[11pt]{article}
\sloppy
\usepackage{xcolor}
\usepackage{scrlayer-scrpage}

\newcommand{\MyName}{P. Lerdkijrachapong}

\usepackage{blindtext}

\clearscrheadfoot
\rohead{\MyName}
\lehead{\today}

\chead{Week 7: Post 1997}
\lohead{02 September 2025}
\rehead{\MyName}    
\cefoot{\thepage}


\begin{document}
Thailand's governance is deeply tied with the democratic system; as such every one has to obey to the Constitution. Since the revolution of 2475 by Khana Rasadorn, Thailand has change its constitution 20 times. With the prominent turning point being its 16th Constitution; the 1997 Constitution, or as we know as the People's Constitution.

The People's Constitution is the first Constitution that successfully integrated the Constitutionalism, which focuses on the people's value. This turning point acts as the foundation of what newer Constitution aims for.

The concept of Constitution is to utilise a power of sovereignty; by splitting the power into 3 main institutions: (1)Legislative, (2)Executive, (3)Judicial. These branches/institutions are what allows us to govern ourselves. But the way that each country use the Constitution will varies according to their form of governance.

Unlike many of the Presidency countries, the usage of parliamentary system has always been implement since the beginning of Thailand's democracy governance. The implementation of Thailand's parliamentary system can be easily explain: Member of Parliament with the Executive power, and Senate with the Legislative power. However, during the beginning period, the concept of formally forming party hasn't been use in Thailand; and will not be use until the end of Semi-Democracy period.


\end{document}
